\documentclass[10pt]{article}
\usepackage[margin=20mm]{geometry}
\usepackage{amsmath,amssymb,mathtools,fontspec,unicode-math,xurl,hyperref}
\IfFontExistsTF{DejaVu Serif}{\setmainfont{DejaVu Serif}}{\setmainfont{Latin Modern Roman}}
\setmathfont{latinmodern-math.otf}
\setlength{\parindent}{0pt}\setlength{\parskip}{6pt}\sloppy
\title{Constraint and boundary checks for singular initial data}
\author{AI technical review}
\date{}
\begin{document}
\maketitle
Abstract: This review examines the necessity argument in Theorem 1.49, concentrating on the critical structure-coefficient condition, the reference-volume change and the limiting momentum constraint. The equality case of the structural condition is supported by a positive uniform decay margin. The apparent logarithmic divergences in the momentum constraint cancel under the stated Hamiltonian identity. No substantive issue was confirmed in this scope. This is a limited initial review, not a verification of all asymptotic estimates or of the proposed global parametrisation of solutions.\par
\section*{Identity and inspected scope}
The paper is Hans Ringström, Initial data on big bang singularities, M25-0272, DOI https://doi.org/10.4171/JEMS/1645. The reviewed file is the EMS Press version of record, J. Eur. Math. Soc. 28 (2026), no. 7, 3049–3130, 82 PDF pages [paper]. Its first page records receipt on 22 August 2022 and revision on 26 September 2024. The publisher separately records acceptance on 26 March 2025 and publication on 9 June 2025 [publisher]. The frozen 2025 membership is preserved; it is distinct from the 2026 issue year. One-based PDF page \(j\) is printed page \(3048+j\).\par
Actual text coverage comprised PDF pp. 1–6, 13–14, 23–29, 59–61, 64–66 and 70–77. The review followed the necessity proof and its momentum constraint calculation, inspected the convergence hypotheses and selected derivative estimates, and performed the independent checks below. All 82 pages were programmatically parsed; that fact does not enlarge the scientific coverage. Sections 2–4, most of Sections 5–6, the full anisotropic bootstrap and the external results cited as [67] and [69] in the paper were not independently proved here. In particular, the curvature identity used before (7.43) remains an explicitly identified dependency.\par
\section*{Hypotheses and logical direction}
Definition 1.10 (PDF p. 13) asks for a Riemannian metric \(\bar h\), a symmetric endomorphism \(\mathfrak K\) with distinct eigenvalues \(p_A\), and scalar data \(\Phi_1,\Phi_0\), satisfying \[\operatorname{tr}\mathfrak K=1,\qquad
 \operatorname{tr}\mathfrak K^2+\Phi_1^2=1,\qquad
 \operatorname{div}_{\bar h}\mathfrak K=\Phi_1\,d\Phi_0.\] For distinct indices it also requires the corresponding eigenframe structure coefficient to vanish on a neighbourhood whenever \(q_{ABC}=1+p_A-p_B-p_C\leq0\) at the point in question.\par
Theorem 1.49 is conditional. Its first part assumes the 2-standard assumptions, the scalar-field estimate (1.32) with \(k=1\), and the limiting Hamiltonian identity \(\bar\Phi_a^2+\operatorname{tr}\mathfrak K^2=1\). The full conclusion adds \(k\)-standard assumptions for \(k\geq2\) and the scalar estimate at that order. The Hamiltonian identity is therefore an input to this necessity theorem, not a conclusion derived from convergence alone.\par
Definitions 1.39–1.45 (PDF pp. 23–25) include a closed spatial manifold, zero shift, positive lapse, a positive lower mean-curvature bound (with the extra stated bound when \(\Lambda>0\)), uniform \(\varrho\to-\infty\), a uniform spectral gap, and spatially differentiated exponential decay with \(\epsilon(x)\geq\epsilon_0>0\). The higher-order lapse assumption controls spatial derivatives of \(\log\widehat N\) polynomially in \(\varrho\). These qualifications cannot be dropped in reading the result.\par
The integration step in Lemma 7.1 is consistent with these quantifiers: \(\widehat U\varrho=1\), and zero shift keeps \(x\) fixed along each integral curve. Thus a bound \(C\langle s\rangle^a e^{2\epsilon(x)s}\) is integrable at \(s=-\infty\), uniformly using \(\epsilon_0\). The claimed limit is not obtained by integrating an arbitrary pointwise convergence statement.\par
\section*{The critical structural condition}
The potentially delicate part of Theorem 1.49 is the neighbourhood conclusion at \(q_{ABC}=0\). The following check reproduces the decisive implication on PDF p. 76, including the direction of the inequality when \(\varrho<0\).\par
Equation reference: eq:positive. Write \([X_B,X_C]=\lambda^A_{BC}X_A\). From (7.31) and (7.43), \[\begin{aligned}
 2\mu_A-2\mu_B-2\mu_C&=2q_{ABC}\varrho+O(1),\\
 \sum_{A,B,C}e^{2\mu_A-2\mu_B-2\mu_C}
       (\lambda^A_{BC})^2&\leq C e^{2\eta\varrho},
       \qquad \eta>0. \end{aligned}\] The error is uniformly bounded near the singularity. If \(q_{ABC}(x_0)\leq0\), continuity supplies a neighbourhood \(V\) on which \(q_{ABC}\leq\eta/2\). For sufficiently negative \(\varrho\), \[2q_{ABC}\varrho+O(1)\geq\eta\varrho-C.\] Every summand in [eq:positive] (\#eq:positive) is nonnegative. Consequently, \[(\lambda^A_{BC})^2\leq C e^{\eta\varrho}\longrightarrow0
 \quad\hbox{on }V.\] The \(C^1\) frame convergence gives convergence of the structure coefficients to those of the limiting frame. Hence the limiting coefficient vanishes throughout \(V\). Rescaling the frame does not alter this vanishing condition for distinct \(A,B,C\).\par
This also disposes of the suggested boundary objection: the proof does not merely infer a neighbourhood statement from pointwise vanishing. The uniform positive \(\eta\) covers nearby small positive values of \(q\). Without a positive rate margin this particular argument would not work. The derivation of [eq:positive] (\#eq:positive) uses the scalar-curvature identity cited in the source and the preceding asymptotic estimates; the present check verifies the implication once those estimates hold and does not replace their full verification.\par
\section*{Reference volume and momentum constraint}
The reference-volume change on PDF p. 28 has the correct sign: \[\bar\varrho=\varrho-r_\theta,\qquad
 \bar\Phi_a=\Phi_1,\qquad
 \bar\Phi_b=\Phi_0+r_\theta\Phi_1.\] Indeed \(\bar\Phi_a\bar\varrho+\bar\Phi_b
=\Phi_1\varrho+\Phi_0\) exactly. Ignoring this change would give the wrong constant term for the scalar field.\par
The cancellation in Lemma 5.21 (PDF pp. 59–61), reused on pp. 76–77, can be checked without interchanging separate divergent limits. For a limiting frame derivative \(D\), the potentially singular combination is \[\sum_A p_A D(-p_A\log\theta)
 +\Phi_1D(-\Phi_1\log\theta+\Phi_0)+D\log\theta .\] Its coefficient of \(\log\theta\) is \(-\sum_Ap_A Dp_A-\Phi_1D\Phi_1\); its coefficient of \(D\log\theta\) is \(1-\sum_Ap_A^2-\Phi_1^2\). Both vanish by the Hamiltonian identity and its differential. The remaining term is \(\Phi_1D\Phi_0\). The exponentially decaying remainders multiplied by polynomial factors still tend to zero.\par
For an orthonormal eigenframe \(E_A\) and \([E_A,E_C]=\gamma^D_{AC}E_D\), the connection terms give \[(\operatorname{div}_{\bar h}\mathfrak K)(E_C)
 =E_Cp_C+\sum_A(p_C-p_A)\gamma^A_{AC}.\] This follows directly by expanding \(\sum_A\bar h((\nabla_{E_A}\mathfrak K)E_C,E_A)\) and using metric compatibility. It agrees with the source’s limiting expression. The rescaling in the proof of Theorem 1.49 places the calculation in the orthonormal frame for the extracted \(\bar h\).\par
\section*{Independent exact examples and limits}
On \(\mathbb T^3\times(0,t_+)\), take \(\Lambda=0\) and \[g=-dt^2+\sum_{i=1}^3t^{2p_i}(dx^i)^2,\qquad
 \phi=a\log t+b,\qquad \sum_i p_i=1,\quad \sum_i p_i^2+a^2=1.\] The nonzero Ricci coefficients are \(R_{00}=(\sum p_i-\sum p_i^2)t^{-2}\) and \(R_{ii}=p_i(\sum p_j-1)t^{2p_i-2}\). Thus \(R_{\alpha\beta}=\partial_\alpha\phi\,\partial_\beta\phi\); the wave equation has coefficient \(a(\sum p_i-1)=0\). Here \(\theta=t^{-1}\), \(\varrho=\log t\), \(\mathcal K=\operatorname{diag}(p_i)\) and \(\widehat U\phi=a\). All rate residuals vanish. The normalized lapse is spatially constant, \(\widehat N=t^{-1}\), so every spatial derivative of \(\log\widehat N\) vanishes. The structure coefficients and both sides of the momentum constraint are zero. The lower bound on \(\theta\) holds on any sufficiently short interval.\par
Exact rational substitution was run for \((p_1,p_2,p_3;a)=(-2/7,3/7,6/7;0)\) and \((0,1/3,2/3;2/3)\). Each has distinct eigenvalues; the corresponding \(q_{123}\) is \(-4/7\) and \(0\), respectively. Every computed Einstein and wave coefficient residual was exactly zero. These are admissible consistency checks, including a critical-wall example; they do not prove the general theorem.\par
No confirmed Major or Minor issue follows from these checks. A bounded title-plus-erratum search and the publisher page did not identify a correction; this is not proof that no correction exists. The supplied manifest records bounded supplement/code searches with no associated material identified. No author code was executed, and absence of optional code is not a mathematical defect.\par
\section*{Conclusion}
Item | Assessment\par
Inspected necessity steps | No Issue Found within the stated scope.\par
Repair | None proposed; no supported defect requiring repair.\par
Decision and adoption | Scientific decision pending; no author adoption requested.\par
Remaining work | Independent curvature/bootstrap verification, remaining sections, and separate acceptance of this initial review.\par
\section*{Supporting checks}
The source PDF has 911510 bytes and SHA256 \[\texttt{0cd01008f97bc714efa9e8ad604f4af741ca86d1f0929da5a76db86f651e0cc1}.\] Batch token: `[internal identifier omitted]`. Its successful command was Python `-B independent\_check.py`; this is coefficient arithmetic, not numerical evolution of Einstein’s equations.\par
A native read-only helper checked the critical-wall implication and rate assumptions; its actual coverage and limitations are recorded separately. The runtime self-description is GPT-6 Astra Pro; the controller’s reported GPT-6/Pro UI selection is a separate observation, not independently verified backend identity. No human review or independent acceptance is claimed. Historical rendering records retain their original false inspection flags.\par
9 H. Ringström, Initial data on big bang singularities, J. Eur. Math. Soc. 28 (2026), 3049–3130, https://doi.org/10.4171/JEMS/1645. Exact acquired publisher PDF. EMS Press, article record for the preceding paper, https://ems.press/journals/jems/articles/14298869, consulted during this review. The supplied manifest also records the associated preprint arXiv:2202.04919; no equivalence comparison was performed.\par
\end{document}
